\scanpage{132}
Par 2 \uline{Etude cohomologique des morphismes \unsure{projectifs}.}

1. \uline{Calculs Explicites de certains groupes de cohomologie}

Soient $X$ un preschema, $\underline L$ un faisceau inversible sur $X$, considerons l'anneau gradue $S=\mathrm I_*(X,\underline L)$ et considerons des $f_i\quad S_{d_i}$ ($1\quad i\quad k$, les $d_i$ des entiers), soit $U_i=X_{f_i}$, $U=\cup U_i$, $\underline U$ le recouvrement \deleted{\unsure{forme}} de $U$ forme par les $U_i$. Si maintenant $F$ est un faisceau quasi-coherent quelconque sur $X$, on posera
\[
\begin{aligned}
F(*)&=\qquad F(\underline L^n\\
H^*(U,F(*))&=\qquad H^*(U,F\quad\underline L^n)\\
H^*(\underline U,F(*))&=\qquad H^*(\underline U,\unsure{R}\quad\underline L^n)
\end{aligned}
\]
Ce sont tous des $S$-modules gradues. Considerons de meme le $S$-module gradue
\[
M=H^0(X,F(*))=\mathrm I(X,F(*))
\]
On aura alors
\[
\mathrm I(U_{i_0i_1\ldots i_p},F(*))=M_{f_{i_0}f_{i_1}\ldots f_{i_p}}=\lim M^{(f_{i_0}f_{i_1}\ldots f_{i_p})^n}
\]
en vertu du Chapitre II, par.1,

en supposant maintenant $X$ noetherien, ou quasi-compact separe); les homomorphismes dans le systeme inductif du dernier membre etant encore la multiplication par $(f_{i_0}f_{i_1}\deleted{\unsure{xxx}}\ldots f_{i_p})^{m-n}$. On en deduit encore, comme au par. precedent:
\[
C^p(\underline U,F(*))=\lim C_n^p(M)
\]
ou $C_n^p(M)$ est forme des fonctions alternees de $p+1$ indices entre $1$ et $k$, a valeurs dans $M$, i.e. pour $p\quad0$ s'identifie a $K^{p+1}(\underline f^n,M)$, les homomorphismes dans ce systeme inductif etant alors les $\underline f^{m-n}$. Par suite on a encore un \deleted{\unsure{homomor}} isomorphisme
\[
C^p(\underline U,F(*))=C^{p+1}\unsure{\}}(\underline f),M)\qquad\text{pour }p\quad0
\]
et ces isomorphismes sont compatibles avec les operateurs differentiels. On fera cependant attention aux questions de degres dans l'isomorphisme precedent; il y a deux degres, celui provenant de $C^*$ et celui provenant de $F(*)$. Dans le deuxieme membre, considere comme limite inductive des $K^{p+1}(\underline f^n,M)$, le \inserted{deuxieme} degre d'une fonction \deleted{\unsure{primitive}} alternee de $p+1$ arguments $i_\alpha$ compris entre $1$ et $k$, a valeurs dans $M_m$, est egal a
\[
m-n(d_{i_1}+d_{i_2}+\ldots+d_{i_p})
\]
Des isomorphismes precedents, on tire la proposition suivante, qui generalise la prop.2 du paragraphe 1:
\margin{131}
